% 第7回　コンセプト立案（記入例・模範解答）
\session{7}{2}{11月11日（水）}{コンセプト立案}{AIと発想法}{yes}

\goals{
  \item 発想法（異業種との掛け合わせ、前提の反転、極端化）をAIとの対話で回せる
  \item 量を出してから選ぶ、という発想の手順を実践できる
  \item コンセプトを「誰に、何を、どのように」の形で言語化できる
}

\afillin{狙うペルソナ（要約）}{A：子連れで自分のペースで比べたい32歳／B：営業トークが苦手な初めての購入者24歳}

\work[60mm]{ワーク1}{AIで30案を出す}
\begin{promptbox}[異業種と掛け合わせる]
自動車ディーラーの店舗に、図書館、ホテルのラウンジ、カフェ、保育園、スポーツジム、公園の要素を掛け合わせた店舗コンセプトを、それぞれ2案ずつ出してください。各案は「誰に／何を／どのように」の形式で、1案3行以内にしてください。ペルソナ：（確定したペルソナの要約）
\end{promptbox}
\begin{promptbox}[前提を反転する]
自動車ディーラーの店舗について「当たり前」とされていることを10個挙げ、それぞれを反転させた場合にどんな店舗になるかを1行ずつ示してください。
\end{promptbox}
\instr{極端化（「もし〜しかなかったら」「もし10倍だったら」）も自分たちで指示を考えて試す。}
\begin{wstabh}{colspec={Q[c,wd=22mm,m] Q[c,wd=12mm,m] X[l,m]},row{2-Z}={ht=9mm}}
  方法 & 案の数 & 気になった案（一言で） \\
  掛け合わせ & \ans{12} & \af{図書館×ディーラー（静かに車の本や資料を読める）／公園×ディーラー（芝生の中に展示車）} \\
  反転 & \ans{10} & \af{「スタッフが話しかける」→「客が呼ぶまで話しかけない」／「商談はテーブル」→「車の中で相談」} \\
  極端化 & \ans{8} & \af{「もし滞在が3時間だったら」→ 子どもが昼寝できる店／「もし展示車が1台だけなら」→ 体験に特化} \\
\end{wstabh}
\afillin{極端化で使った自分たちの指示}{「来店客が3時間いても飽きない店にするとしたら、何が必要か。5案」}

\work[70mm]{ワーク2}{5案に絞る（ペルソナの言葉で判定）}
\begin{wstabh}{colspec={Q[c,wd=6mm,m] X[2,l,m] X[3,l,m] Q[c,wd=14mm,m]},row{2-Z}={ht=13mm}}
  & 案 & ペルソナは「行きたい」と言うか（本人の言葉で） & 判定\newline ○△× \\
  1 & \af{呼ぶまで話しかけない店（サインカード式）} & \af{B「これなら、気軽に見に行ける」} & \ans{○} \\
  2 & \af{公園のようなキッズスペース併設} & \af{A「子どもが遊んでいる間に、ゆっくり見られる」} & \ans{○} \\
  3 & \af{図書館のような資料ラウンジ} & \af{B「調べた内容を確かめられる」 A「子連れだと静かすぎて気を使う」} & \ans{△} \\
  4 & \af{車の中で相談できる商談スペース} & \af{B「面白い。でも本当に落ち着くかは不安」} & \ans{△} \\
  5 & \af{ジム併設で整備待ちに運動} & \af{A・B「行く理由にはならない」} & \ans{×} \\
\end{wstabh}

\work[70mm]{ワーク3}{コンセプト文を書く}
\begin{wstab}{colspec={Q[l,wd=30mm,m] X[l,m]},row{1-Z}={ht=13mm}}
  \hc{誰に} & \af{子育て世帯と、初めて車を買う若い世代に} \\
  \hc{何を} & \af{自分のペースで車を比べ、聞きたいときだけ相談できる時間を} \\
  \hc{どのように} & \af{呼ぶまで話しかけないサインカードと、子どもを見守れるキッズスペースのある、公園のようなラウンジで} \\
  \hc{コンセプト文\newline（一文で）} & \af{「見るだけ、大歓迎。」子どもも自分も、自分のペースで過ごせる“公園ラウンジ”のディーラー} \\
  \hc{予備案} & \af{調べてきたことを確かめられる「答え合わせ」の店（資料ラウンジ＋試乗の予約制）} \\
\end{wstab}

\work{発表}{各グループ1分。他グループの案で印象に残ったもの}
\alines{3}{グループA「夜だけ開く大人のガレージ」：時間帯で客層を分ける発想が新しい。\par
グループD「整備待ちの時間を体験に変える」：既存の客（整備）を入口にする点は、自分たちの案にも取り入れたい。}

\begin{nexttime}
  \item コンセプトシート1枚（コンセプト文、狙うペルソナ、根拠、参考にした事例）
  \item コンセプトを実現する空間に必要な要素を、思いつく限り書き出す
\end{nexttime}
\atwoboxes{根拠（企業の課題・ペルソナとの関係）}{参考にした事例（出典）}{32mm}
{課題は「若い世代や子育て世帯の来店を増やす」こと。ペルソナA・Bに共通する不安は「話しかけられて自分のペースで見られない」ことで、ここを解消することが来店の動機になる。}
{子連れ向けのスペースを設けた商業施設、セルフ方式の家具店の売り場など（各自が訪問・調査した事例を、店名・確認日とともに記入する）}
\abox[空間に必要な要素（思いつく限り）]{28mm}{サインカード（「見ているだけ」「相談したい」）、入口から見えるキッズスペース、芝生風の床、自由に座れるソファ、車の資料と価格の目安を置いた棚、セルフで乗り込める展示車、ガラス越しに見える整備工場、授乳室、カフェカウンター}
\aimemoA{店舗コンセプトの発想（掛け合わせ・反転・極端化）}
{「異業種6つと掛け合わせて2案ずつ」「当たり前を10個挙げて反転」「3時間いても飽きない店にするには」}
{30案から5案に絞り、ペルソナの言葉で判定して1案に決めた。コンセプト文はAIの案を参考にせず、グループで書いた。}
{「呼ぶまで話しかけない」接客が実際にある業種を調べ、事例を確認した}

\begin{point}
  \item 「量を出してから選ぶ」手順が守られているか（30案→5案→1案）を見る。最初から1案に決めているグループには、選ばなかった案を記録させる。
  \item 判定は、自分たちの好みではなくペルソナの言葉で行っているかが大切。
  \item コンセプト文は「誰に・何を・どのように」の3要素がそろい、第6回の課題と結びついているかで評価する。キャッチコピーの出来は第11回で扱うので、ここでは問わない。
  \item 参考事例は、店名と確認日を書かせる。AIが挙げた事例は実在と内容を確かめさせる。
\end{point}
